\documentclass[12pt,a4paper]{report}

\usepackage{graphicx}
\usepackage{amsmath}
\usepackage{booktabs}
\usepackage{array}
\usepackage{float}

% ==================================================
% TITLE PAGE
% ==================================================

\title{StitchPerfect}
\author{\texttt{1024160031} Uday Singh}
\date{Project Evaluation\\[0.2cm]
Computer Engineering Department\\[0.2cm]
Thapar Institute of Engineering and Technology, Patiala}

% ==================================================

\begin{document}

\maketitle

\tableofcontents

% ==================================================
% CHAPTER 1
% ==================================================

\chapter{Introduction}

\section{Background and Context}

Tailoring businesses handle a large amount of customer information, including
customer details, body measurements, clothing orders, and measurement history.
When these records are maintained manually, it can become difficult to retrieve
old measurements, update customer information, track orders, and maintain
accurate records.

The proposed StitchPerfect provides a
centralized software solution for managing these activities. The system allows
the administrator to store and manage customer information, body measurements,
and order records through a graphical user interface.

An additional feature of the system is an automatic measurement prediction
module. Given basic inputs such as height and weight, the system uses a
measurement dataset containing approximately 10,000 records to estimate
multiple body measurements. This can assist the tailor in obtaining an initial
set of measurements before manually verifying them.

\section{Problem Statement}

Traditional tailoring workflows often depend on handwritten registers or
separate records. Such approaches can result in duplicated information,
difficulty in searching previous records, errors while copying measurements,
and inefficient order management.

The project addresses these problems by developing an integrated application
that stores customer and tailoring records in a database and provides
functionality for adding, viewing, updating, and managing these records.
Furthermore, the system introduces an automated prediction feature that uses
height and weight as inputs to estimate body measurements.

\section{Project Objectives}

The main objectives of the project are:

\begin{enumerate}
  \item To develop a centralized system for managing customer information.
  \item To store and manage body measurements digitally.
  \item To maintain tailoring order records in a structured database.
  \item To provide CRUD functionality for the major system entities.
  \item To develop an automatic body-measurement prediction feature using
        height and weight.
  \item To use a measurement dataset of approximately 10,000 rows for the
        prediction component.
  \item To provide a simple graphical interface that can be used efficiently
        by a tailor or administrator.
\end{enumerate}

\section{Scope of the Project}

The system focuses on the management of customers, measurements, and orders
within a tailoring environment. The administrator is the primary user who
enters and manages records.

The prediction component is intended to provide estimated measurements from
basic physical inputs. The predicted values can be reviewed and verified by
the tailor before being used for actual garment preparation.

% ==================================================
% CHAPTER 2
% ==================================================

\chapter{Methodology and Design}

\section{System Architecture}

The proposed system follows a layered application structure consisting of a
graphical user interface, application logic, prediction functionality, and a
database layer.

The administrator interacts with the web interface. The application
processes the entered information and performs the required database
operations. Customer details, measurements, and order information are stored
in MySQL. The prediction module receives height and weight and produces
estimated measurement values using the prepared measurement dataset.

\subsection{High-Level Design}

The major components of the system are:

\begin{itemize}
  \item \textbf{Admin Interface:} Provides access to the management
        functionality.
  \item \textbf{Customer Management:} Handles customer records and their
        associated information.
  \item \textbf{Measurement Management:} Stores and retrieves body
        measurements.
  \item \textbf{Order Management:} Maintains tailoring order information.
  \item \textbf{Measurement Predictor:} Estimates body measurements from
        height and weight.
  \item \textbf{Database Layer:} Stores persistent customer, measurement,
        and order records using MySQL.
\end{itemize}

\subsection{Technical Stack}

\begin{itemize}
  \item \textbf{Frontend:} HTML, CSS, JavaScript and Bootstrap
  \item \textbf{Backend:} Node.js and Express.js
  \item \textbf{Database:} MySQL
  \item \textbf{Database Connectivity:} MySQL driver for Node.js
  \item \textbf{Prediction Component:} Body-measurement prediction model
        based on a prepared measurement dataset
  \item \textbf{Development Tools:} Visual Studio Code and Git
\end{itemize}

\section{Database Design}

The database provides persistent storage for the system. The major logical
entities are Customer, Measurement, and Order.

\subsection{Customer Entity}

The Customer entity stores information required to identify and manage a
customer. A customer can have measurement records and multiple order records
over time.

\subsection{Measurement Entity}

The Measurement entity stores body measurements required for tailoring. The
system can maintain measurements such as:

\begin{itemize}
  \item Chest
  \item Waist
  \item Hips
  \item Shoulder
  \item Sleeve
  \item Inseam
  \item Neck
\end{itemize}

Additional measurement attributes can be included according to the final
database design of the project.

\subsection{Order Entity}

The Order entity stores information related to garments ordered by a
customer. It can be used to track the order and associate it with the
corresponding customer and measurements.

\section{Implementation Details}

\subsection{Customer Management Module}

The customer management module allows the administrator to enter customer
details and maintain existing records. CRUD operations are used to create,
retrieve, update, and delete records as required.

\subsection{Measurement Management Module}

The measurement management module provides facilities for recording and
retrieving customer measurements. Measurements can be associated with the
corresponding customer so that previous records can be accessed when
required.

\subsection{Order Management Module}

The order management module maintains tailoring orders and their associated
customer information. This helps the administrator keep track of the work
that has been requested for each customer.

\subsection{Automatic Measurement Prediction}

The prediction module is designed to estimate body measurements using basic
inputs, particularly height and weight. The system uses a prepared dataset
containing approximately 10,000 measurement records.

The prediction workflow is:

\begin{enumerate}
  \item The administrator enters the customer's height and weight.
  \item The prediction component processes these input values.
  \item The trained or prepared prediction logic estimates the required body
        measurements.
  \item The predicted measurements are displayed to the administrator.
  \item The administrator can verify the values before using them for
        tailoring.
\end{enumerate}

The prediction feature is intended as an assistance mechanism rather than a
replacement for physical measurement and professional verification.

\section{System Workflow}

The general workflow of the application is:

\begin{enumerate}
  \item Administrator opens the application.
  \item Customer information is entered or an existing customer is selected.
  \item Measurements can be entered manually or estimated using the
        prediction feature.
  \item Measurement records are stored in the database.
  \item An order is created and associated with the customer.
  \item The stored information can later be retrieved and updated.
\end{enumerate}

% ==================================================
% CHAPTER 3
% ==================================================

\chapter{Results and Evaluation}

\section{Testing Methodology}

The system can be evaluated through functional and integration testing of its
major modules.

\begin{itemize}
  \item \textbf{Customer Testing:} Verify creation, retrieval, modification,
        and deletion of customer records.
  \item \textbf{Measurement Testing:} Verify that measurement records are
        stored and retrieved correctly.
  \item \textbf{Order Testing:} Verify that orders are correctly associated
        with customers.
  \item \textbf{Database Testing:} Verify successful insertion, retrieval,
        updating, and deletion of database records.
  \item \textbf{Prediction Testing:} Verify that height and weight inputs
        produce measurement estimates.
  \item \textbf{Integration Testing:} Verify communication between the
        web interface, application logic, prediction component, and
        MySQL database.
\end{itemize}

\section{Functional Results}

The completed system is designed to provide a single interface through which
the administrator can manage the major tailoring activities.

\begin{table}[H]
\centering
\begin{tabular}{|p{4cm}|p{8cm}|}
  \hline
  \textbf{Module} & \textbf{Expected Result} \\
  \hline
  Customer Management & Customer records can be added, viewed, updated,
  and deleted. \\
  \hline
  Measurement Management & Body measurements can be stored and retrieved
  for customers. \\
  \hline
  Order Management & Tailoring orders can be created and associated with
  customers. \\
  \hline
  Prediction Module & Height and weight can be used to generate estimated
  body measurements. \\
  \hline
  Database & Records remain available after being stored in MySQL. \\
  \hline
\end{tabular}
\caption{Functional Results Summary}
\label{tab:functional-results}
\end{table}

\section{Prediction Dataset}

The prediction component uses a body-measurement dataset containing
approximately 10,000 rows. The dataset provides reference examples for
estimating body measurements from input characteristics.

The final report should include the exact prediction-model metrics, dataset
statistics, and validation results obtained during implementation. These
values should be inserted after the final model has been evaluated.

\section{Analysis}

The system combines traditional record management with an automated
prediction feature. Centralized database storage improves the organization
and retrieval of customer information, while the prediction component can
reduce the effort required to obtain an initial measurement estimate.

The predicted measurements should always be reviewed by the tailor because
body measurements can vary between individuals even when their height and
weight are similar.

% ==================================================
% CHAPTER 4
% ==================================================

\chapter{Conclusion}

\section{Summary of Work}

The StitchPerfect provides a digital
solution for managing customer information, measurements, and tailoring
orders. The application uses a web interface and MySQL database to provide
structured storage and management of records.

The project also incorporates an automated measurement prediction component
that uses height and weight as inputs and a large body-measurement dataset as
the reference for generating estimated measurements.

\section{Key Findings}

\begin{itemize}
  \item A centralized database can simplify the management of tailoring
        records.
  \item Digital measurement records make previously stored customer
        information easier to retrieve and update.
  \item Automated prediction can provide an initial estimate of body
        measurements from basic inputs.
  \item Prediction results require verification before being used for actual
        garment production.
\end{itemize}

\section{Future Work and Recommendations}

Possible future improvements include:

\begin{enumerate}
  \item Improving the prediction model by using a larger and more diverse
        measurement dataset.
  \item Adding additional input features to improve measurement estimates.
  \item Adding authentication and role-based access control.
  \item Adding customer order status and notification functionality.
  \item Providing invoice and bill generation.
  \item Adding measurement history and comparison features.
  \item Developing a web or mobile version of the system.
  \item Evaluating multiple machine-learning algorithms and comparing their
        performance using suitable metrics.
\end{enumerate}

\section{Final Remarks}

The project demonstrates how software engineering, database management,
graphical user interfaces, and predictive techniques can be combined to
digitize a traditional tailoring workflow. The system provides a foundation
for further automation and can be extended with additional management and
prediction capabilities.

% ==================================================
% BIBLIOGRAPHY
% ==================================================



\end{document}
